
幺正规范把三个 Goldstone 方向完全吸收到有质量规范场中，因此最适合读取树级质量谱；圈图计算则需要同时控制纵向规范场、Goldstone 场与鬼场的紫外行为。$R_\xi$ 规范通过一个连续参数 $\xi_V$ 把这些非物理自由度放在相同的规范依赖极点上，从而使传播子的高动量幂次计数保持透明。

为保持 SI 量纲一致，把具有 $[\Phi]=\sqrt{\mathrm{J/m}}$ 的标量涨落除以 $\sqrt{\hbar c}$，定义几何量纲场
\begin{equation*}
 \widehat\phi^\pm\equiv\frac{\phi^\pm}{\sqrt{\hbar c}},
 \qquad
 \widehat\chi\equiv\frac{\chi}{\sqrt{\hbar c}},
 \qquad
 \widehat h\equiv\frac{h}{\sqrt{\hbar c}},
\end{equation*}
它们的 SI 单位均为 $\mathrm{m^{-1}}$。同时定义 $W,Z,h$ 的 Compton 逆长度
\begin{equation*}
 \kappa_W\equiv\frac{m_Wc}{\hbar},
 \qquad
 \kappa_Z\equiv\frac{m_Zc}{\hbar},
 \qquad
 \kappa_h\equiv\frac{m_hc}{\hbar}.
\end{equation*}
在笛卡尔参数化中，Higgs 二重态写成
\begin{equation}
 \Phi=
 \begin{pmatrix}
 \phi^+\\[2pt]
 \dfrac{v+h+\ii\chi}{\sqrt2}
 \end{pmatrix},
 \qquad \phi^-=(\phi^+)^*.
 \label{eq:sm-higgs-cartesian-rxi-dp11}
\end{equation}
把它代入 $|D_\mu\Phi|^2-V(\Phi)$ 并只保留真空附近的二次项，得到规范场与标量涨落尚未解耦的二次拉氏量
\begin{align}
 \frac{\Lag_{\Phi}^{(2)}}{\hbar c}
 ={}&\partial_\mu\widehat\phi^-\partial^\mu\widehat\phi^+
 +\frac12(\partial_\mu\widehat h)^2
 +\frac12(\partial_\mu\widehat\chi)^2
 -\frac12\kappa_h^2\widehat h^2
 \nonumber\\
 &+\kappa_W^2\mathcal W_\mu^+\mathcal W^{-\mu}
 +\frac12\kappa_Z^2\mathcal Z_\mu\mathcal Z^\mu
 \nonumber\\
 &+\ii\kappa_W
 \left(
 \mathcal W^-_\mu\partial^\mu\widehat\phi^+
 -\mathcal W^+_\mu\partial^\mu\widehat\phi^-
 \right)
 +\kappa_Z\mathcal Z_\mu\partial^\mu\widehat\chi.
 \label{eq:sm-rxi-quadratic-before-gf-dp11}
\end{align}
最后一行说明直接求逆二次核时，$W_\mu^\pm$ 与 $\phi^\pm$、$Z_\mu$ 与 $\chi$ 会发生双线性混合。$R_\xi$ 规范固定函数取为
\begin{align}
 F^+&=\partial^\mu\mathcal W_\mu^+
 -\ii\xi_W\kappa_W\widehat\phi^+,
 &F^-&=\partial^\mu\mathcal W_\mu^-
 +\ii\xi_W\kappa_W\widehat\phi^-,\\
 F^Z&=\partial^\mu\mathcal Z_\mu
 -\xi_Z\kappa_Z\widehat\chi,
 &F^\gamma&=\partial^\mu A_\mu^{\rm geo},
 \label{eq:sm-rxi-gauge-functions-dp11}
\end{align}
并加入
\begin{equation}
 \boxed{
 \Lag_{\rm gf}
 =-\frac{\hbar c}{\xi_W}F^+F^-
 -\frac{\hbar c}{2\xi_Z}(F^Z)^2
 -\frac{\hbar c}{2\xi_\gamma}(F^\gamma)^2.}
 \label{eq:sm-rxi-gf-lagrangian-dp11}
\end{equation}
这些函数中的标量项并非任意附加：其系数恰好由\eqnrefs{eq:sm-rxi-quadratic-before-gf-dp11} 的规范场--Goldstone 混合项固定。展开并分部积分后，所有双线性混合精确相消，同时 Goldstone 场获得规范依赖质量
\begin{equation}
 \boxed{
 m_{\phi^\pm}^2=\xi_Wm_W^2,
 \qquad
 m_{\chi}^2=\xi_Zm_Z^2.}
 \label{eq:sm-rxi-goldstone-masses-dp11}
\end{equation}
这里的质量只决定规范固定后的内部传播子极点，不代表新的物理粒子谱线。采用物理四动量 $p^\mu$ 后，相应传播子为
\begin{equation*}
 \Delta_{\phi^\pm}(p)
 =\frac{\ii}{p^2-\xi_Wm_W^2c^2+\ii0^+},
 \qquad
 \Delta_{\chi}(p)
 =\frac{\ii}{p^2-\xi_Zm_Z^2c^2+\ii0^+}.
\end{equation*}

规范固定以后还必须修正路径积分中对同一规范轨道的重复计数。设 $\alpha^b(y)$ 是无穷小规范变换参数，$F^a(x)$ 是所选规范条件。把规范条件沿规范轨道改变多少的线性响应写成
\begin{equation*}
 \mathcal M_{\rm FP}^{ab}(x,y)
 \equiv
 \frac{\delta F^a(x)}{\delta\alpha^b(y)}.
\end{equation*}
这个线性算符称为 Faddeev--Popov 算符。它的行列式正是规范轨道与规范固定截面之间的 Jacobian；在量子场论中通常用一对反交换标量场 $c^a,\bar c^a$ 表示这个行列式，这些辅助场称为 Faddeev--Popov 鬼场。鬼场不是外部可探测粒子，它们只在内部传播子和圈图中补偿非物理规范自由度。

在线性化真空附近，破缺生成元使 Goldstone 场沿规范轨道平移，因此 $\mathcal M_{\rm FP}$ 的二次部分与 Goldstone 二次核具有相同的规范依赖质量尺度。由此得到鬼场质量
\begin{equation}
 \boxed{
 m_{c^\pm}^2=\xi_Wm_W^2,
 \qquad
 m_{c^Z}^2=\xi_Zm_Z^2,
 \qquad
 m_{c^\gamma}=0.}
 \label{eq:sm-rxi-ghost-masses-dp11}
\end{equation}
与 QED 不同，电弱 Faddeev--Popov 行列式依赖 Higgs 场和非阿贝尔规范场，所以电弱鬼场在圈图中不脱耦。

对任一有质量规范玻色子 $V=W,Z$，$R_{\xi_V}$ 规范下的约化传播子统一为
\begin{equation}
 \boxed{
 D^{(V)}_{\mu\nu}(p)
 =\frac{-\ii}{p^2-m_V^2c^2+\ii0^+}
 \left[
 \eta_{\mu\nu}
 -(1-\xi_V)
 \frac{p_\mu p_\nu}{p^2-\xi_Vm_V^2c^2+\ii0^+}
 \right].}
 \label{eq:sm-rxi-massive-vector-propagator-dp41}
\end{equation}
横向物理极点位于 $p^2=m_V^2c^2$，纵向规范极点位于 $p^2=\xi_Vm_V^2c^2$，后者与相应 Goldstone 和鬼场极点重合。$\xi_V=1$ 给出 Feynman--'t Hooft 规范
\begin{equation*}
 D_{\mu\nu}^{(V)}(p)
 =\frac{-\ii\eta_{\mu\nu}}{p^2-m_V^2c^2+\ii0^+},
\end{equation*}
而 $\xi_V\to\infty$ 把非物理极点推向无穷，并得到幺正规范传播子
\begin{equation*}
 D_{\mu\nu}^{(V)}(p)
 \longrightarrow
 \frac{-\ii}{p^2-m_V^2c^2+\ii0^+}
 \left(\eta_{\mu\nu}-\frac{p_\mu p_\nu}{m_V^2c^2}\right).
\end{equation*}
幺正规范的纵向分子在 $|p^2|\gg m_V^2c^2$ 时不随动量衰减，因此逐图紫外幂次计数较差；有限 $\xi_V$ 的 $R_\xi$ 规范把可重整化结构显式保留在传播子中。物理 $S$ 矩阵最终与 $\xi_V$ 无关，这一结论由 BRST 与 Slavnov--Taylor 恒等式保证，而不是要求每一张 Feynman 图分别与 $\xi_V$ 无关。

\begin{figure}[htbp]
\centering
\begin{tikzpicture}[
  >=Latex,
  node distance=1.15cm and 1.05cm,
  box/.style={draw,rounded corners,align=center,text width=3.45cm,minimum height=1.25cm,inner sep=4pt},
  line label/.style={fill=white,fill opacity=.96,text opacity=1,inner sep=1.0pt,font=\footnotesize}
]
\node[box] (mix) {quadratic EWSB kernel\\gauge--Goldstone mixing};
\node[box,right=of mix] (gf) {$R_\xi$ gauge fixing\\choose functions $F^a$};
\node[box,right=of gf] (sector) {gauge/Goldstone propagators\\+ FP ghost propagators};
\node[box,below=of gf] (brst) {BRST symmetry\\Slavnov--Taylor identities};
\node[box,below=of sector] (sum) {sum all gauge-related\\diagrams in one amplitude};
\node[box,below=of sum] (obs) {physical amplitude\\independent of $\xi$};
\draw[->] (mix)--node[line label,midway,above]{cancel mixing}(gf);
\draw[->] (gf)--node[line label,midway,above]{fix propagators}(sector);
\draw[->] (gf)--node[line label,midway,left]{gauge-fixed symmetry}(brst);
\draw[->] (sector)--node[line label,midway,right]{all internal sectors}(sum);
\draw[->] (brst)--node[line label,midway,below,sloped]{identity constraints}(sum);
\draw[->] (sum)--node[line label,midway,right]{$\xi$ cancels}(obs);
\end{tikzpicture}
\caption{$R_\xi$ 规范固定先消去二次 gauge--Goldstone 混合，并同时确定纵向规范场、Goldstone 与 Faddeev--Popov ghost 的传播结构。完整振幅必须把这些规范相关内部扇区与物理图一起求和；BRST 对称性导出的 Slavnov--Taylor 恒等式约束各项，最终使可观测振幅与 $\xi$ 无关。}
\label{fig:sm-rxi-pole-structure-dp11}
\end{figure}
